% 原书第 335 页（PDF 第 348 页）
\BookPage{335}{348}
\BookHeading{试析雍正乾隆父子对待宗族的差异}{试析雍正乾隆父子对待宗族的差异}{陈龙贵*}
「敬天法祖」是有清一代的祖制，「三年不改父志」则是儒家的孝道传统，清代帝王因此一向强调「以孝治天下」。雍正崩殂，乾隆继位，照说应该是子承父志；实则乾隆登极后，种种举措，譬如驱逐道士出宫、赶走和尚文觉、不许陈奏祥瑞之事、为宗室皇亲允禩（阿其那）、允禟（塞思黑）等翻案、凌迟处死曾静、张熙、收缴销毁《大义觉迷录》等，不仅没有率由旧章，倒像个不肖子了¹。而在对待宗族的态度上，乾隆与乃父雍正也存在著作法上的差异；这其间产生差异之处，即本文尝试分析所在。

一、雍正对待宗族的态度：扶持、利用宗族力量

雍正对待宗族的态度，首先，我们可以用以下两则资料来加以说明。

\begin{BookQuote}
乙丑(初十日）…刑部议覆。署江西巡抚迈柱奏，永新县民朱伦三、同姪朱三杰，致死伊弟朱宁三一案，朱伦三应拟流徙，朱三杰应拟徒。得旨：从来凶悍之人，偷窃奸宄，怙恶不悛，以致伯叔兄弟重受其累。本人所犯之罪，在国法虽未至于死；而其尊长族人，翦除凶恶，训诫子弟，治以家法，至于身死，亦是惩恶防患之道，情非得已，不当按律拟以抵偿。如朱伦三因伊弟朱宁三，屡次犯窃，累伊鬻男变产、代赔赃银；又复偷牛被获，故将朱宁三致死。朱三杰立未与谋。着将朱伦三、朱三杰徒流等罪，俱从宽免。嗣后凡遇凶恶不法之人，经官惩治，怙恶不悛，为合族所共恶者，准族人鸣之于官。或将伊流徙远方，以除宗族之害；或以家法处治，至于身死，免其抵罪。着九卿详悉定议
\end{BookQuote}

\vfill
\begin{BookNotes}
\BookNote{}{作者任职于国立故宫博物院。}
\BookNote{1.}{参见陈捷先老师，《乾隆写真》（台北：远流出版公司，2005年5月，初版五刷），〈新君乾隆是不肖子?，页39-46。}
\end{BookNotes}

% 原书第 336 页（PDF 第 349 页）
\BookPage{336}{349}
\begin{BookQuote}
具奏。寻议，凶悍之人，伯叔兄弟治以家法，因而致死，若必按律拟抵，则不法子弟，终不知所儆惧。嗣后许族人呈明地方官，照所犯罪科断。若已经官惩治，仍不悛改，该地方官查明过犯实迹，流三千里。傥事起一时，合族公愤，处以家法致死；该地方官审明所犯，确有应死之罪，将为首者照罪人应死而擅杀律予杖；若罪不至死，将为首者照应得之罪减一等，免其抵偿。若本人竝非凶悍不法，无过犯实迹，而族人诬捏殴毙者，将为首之人，仍照本律科断。从
\end{BookQuote}

之。2

这一则刑部议覆的资料显示：（一）因署江西巡抚迈柱的奏报，雍正皇帝主动降谕旨，认为「着将朱伦三、朱三杰徒流等罪，俱从宽免」，免除了朱伦三、朱三杰的罪刑。（二）而且进一步声明「嗣后凡遇凶恶不法之人，经官惩治，怙恶不悛，为合族所共恶者，准族人鸣之于官。或将伊流徙远方，以除宗族之害；或以家法处治，至于身死，免其抵罪」，让「九卿详悉定议具奏」。于是刑部仰承上意，具奏：（一）朱伦三、同姪朱三杰，致死伊弟朱宁三一案，朱伦三、朱三杰原定流徒的罪行，获得了宽免。（二）因此案下旨定议后的具奏结果是:1.「傥事起一时，合族公愤，处以家法致死；该地方官审明所犯，确有应死之罪，将为首者照罪人应死而擅杀律予杖」；2.「若罪不至死，将为首者照应得之罪减一等，免其抵偿」；3.「若本人竝非凶悍不法，无过犯实迹，而族人诬捏殴毙者，将为首之人，仍照本律科断」。

雍正同意了刑部的议覆，这一结果使得家族取得了一部分「处死不法族人」的权力。从所引实录的文字，审酌其意，表面上似乎是雍正被动的同意刑部的议覆，实则是雍正主动的作为。雍正为何要给予家族这种权力呢？

\begin{BookQuote}
四年，严饬力行保甲，定保正甲长牌头赏罚及选立族正之例。…如有堡子村庄，聚族满百人以上，保甲不能徧查者，拣选族中人品刚方，素为阖族敬惮之人，立为族正；如有匪类，报官究治，徇情隐匿者，与保甲一体治罪。3
\end{BookQuote}

为了安靖地方，「如有堡子村庄，聚族满百人以上，保甲不能徧查者」，于族中选立族正；而此一文字记载，也是雍正主动降下谕旨让「九卿详议具奏」的。无论是雍正

\vfill
\begin{BookNotes}
\BookNote{2.}{《大清世宗宪皇帝实录》，卷五十七，雍正五年五月乙丑条。}
\BookNote{3.}{（清）高宗敕撰，《清朝文献通考》，卷二十三，〈职役考三〉，「雍正四年」。}
\end{BookNotes}

% 原书第 337 页（PDF 第 350 页）
\BookPage{337}{350}
\noindent 四年的选立族正，或是雍正五年「朱宁三案」处死不法族人的权力，家族都从中获取了一些公权力，显见得雍正在这方面，对家族是给予正面扶持的；另一方面，雍正所以这么做，则在利用家族的力量，使地方得以安靖。

明清以来，退休官员是构成在籍士绅的中坚成员；而士绅又是地方上组成家族的重要力量。可是我们从侧面了解到，雍正继位之初，退休官员在地方上似乎并不是那么的受到重视：

\begin{BookQuote}
雍正元年。癸卯。春。正月。辛巳朔。…颁发训谕直省总督以下等官上谕十一道。谕总督。…而家居之大臣官员。皆经皇考任用之人。致政投闲。安静自好。为大吏者、并不加以礼貌。或反被土豪侵欺者有之。亦非所以表率属僚、彰善瘅恶之道也。…谕按察司。…至于良善巨室。致政大臣。猾吏或借端侵抑之。非所以称朕意也。4
\end{BookQuote}

为了地方上的安靖，雍正要拉拢致仕大员，甚至拉拢由这些在籍士绅组成的家族力量。然而雍正继位之初时的退休在籍官员与地方家族，从元年元旦的十一道谕旨来看，似乎不是地方上的重要力量，反而是「被土豪侵欺者」、「猾吏或借端侵抑之」；因此雍正一方面在给直省总督与按察使的上谕里，要他们照顾好退休的大臣，一方面也拉拢这些致政大臣背后所组成的家族力量。

「安靖地方」、「地方安静」，历来是帝王政权治理的目标，因为地方安静了，统治才能顺畅。因此，无论是拉拢退休大臣及其背后的家族力量，或是给予家族某些公权力，无非是雍正要扶持、利用家族的力量，来安靖地方。然而雍正扶持、利用家族，仅仅只是安靖地方而已吗？他那么积极的作为，难道没有另外的目的吗？

二、试析雍正扶持、利用宗族力量的原因

以上说明了雍正拉拢退休大员及扶持地方家族的做法，与其安靖地方作用的原因。而扶持利用地方家族只是雍正所做众多地方事务中很小的一部分，因此雍正扶持利用家族力量的原因，还可以放在地方事务工作的脉络中进一步加以分析。

\vfill
\begin{BookNotes}
\BookNote{4.}{《大清世宗宪皇帝实录》，卷三，雍正元年正月辛巳（元旦）条。}
\end{BookNotes}

% 原书第 338 页（PDF 第 351 页）
\BookPage{338}{351}
\BookSubheading{（一）雍正重视地方事务}

雍正继位之前，在皇子与亲王时期，除了随扈康熙、代理祭孔、祭礼等等之外，似乎没有主持地方事务的经验；可是在他继位之后，却做了大量的地方事务工作，举其荦荦大者如下：

\begin{BookQuote}
雍正元年元旦：颁发训谕直省总督以下等官上谕十一道雍正元年十月：摊丁入亩或丁随地起（以直督李维钧始）雍正元年二年：火耗归公（以诺岷、石文焯为代表）雍正二年二月：颁行圣谕广训并御制序雍正四年四月:士民一例当差雍正四年、五年：改土归流（以鄂尔泰的工作为中心）雍正五年四月：豁贱为良(以豁除江南世仆贱籍之上谕为中心)雍正八年三月:为《钦颁州县事宜》颁发上谕
\end{BookQuote}

诚如陈捷先老师所说，雍正很多的作为，都是为了清朝的统制着想的，甚至在「民无二王」的独尊思想下，他自己是天下唯一的主宰。同样的，雍正所做的大量地方事务，是为了统治着想；雍正扶持利用宗族的力量，目的是安靖地方；都是为了清朝的统制能够长治久安，当然更是为了巩固他的君权着想的。

\BookSubheading{（二）雍正重视地方事务的原因分析}

雍正继位之初，他的君权是不稳固的。这要追溯到康熙晚年时期的储位之争。

康熙四十七年首度废储后，不料竟然启动了皇子之间争夺储位的斗争；康熙为了消弭诸皇子之争，其后虽然又立储、再度废储，可是终康熙之世，储位之争毫无止歇之势。这其间，雍亲王（雍正继位前的爵位）采取的是隐而不显的方式；而争储最力

\vfill
\begin{BookNotes}
\BookNote{5.}{陈捷先老师，《雍正—勤政的皇帝·传奇的一生》（台北：国立故宫博物院，民国九十八年九月初版）；陈捷先老师，〈法家天子—雍正皇帝〉（收入《两岸故宫第一届学术研讨会：为君难一雍正其人其事及其时代》，台北：国立故宫博物院，民国九十九年八月初版），页137-167。}
\end{BookNotes}

% 原书第 339 页（PDF 第 352 页）
\BookPage{339}{352}
\noindent 的则是以皇八子允异为中心的集团，他们是最活跃、也是当时最有实力的集团。

雍正继位之初君权的不稳固，主要也来自允禩、允禟、允俄、允的这一集团夺权与不合作。然而雍正在争储过程中，因为隐而不显，实力不强；而允异集团的实力则是盘根错节，自中央以至地方都有他们的支持者。雍正要稳固君权的策略安在？

虽然雍正君权不稳，他既然继位为皇帝，他就有「君位皇权」的「势」。在中央，雍正暂时隐忍，除了分化允禩集团外，给了允禩高位（封为廉亲王）作为安抚；可是对于地方，他运用至高无上的「势」，争取地方督抚等诸侯的支持。于是在雍正元年元旦，距康熙六十一年十一月十三日康熙之死不过四十几天的时间，还在为逝去的康熙皇帝举丧期间、还是百事待举之时，属于他为君后的第一个元旦，雍正不但不举行朝贺的大典，也不稍事休息，反而迫不及待地向地方自督抚以至州县连连颁发了十一道上谕；为的就是向地方官员以及间接的向全国百姓表示，他才是继康熙之位的皇帝。雍正首先向地方争取支持（甚至可以说是与允禩集团在地方上的夺权竞赛）、巩固君权的策略，我们也可以从《大义觉迷录》中看出：

\begin{BookQuote}
…暂将阿其那拘禁，降旨询问外省封疆大臣，待其回奏，然后定夺。8
\end{BookQuote}

试问，外省封疆大臣接到了皇帝的降旨询问，他们能怎么做？能不表态输诚吗？尤其是原来支持允禩的那些大臣们，他们能不诚惶诚恐：这该不会是皇帝来试探我们的吧？当帝王的「势」面对地方封疆大臣时，仍然有着无比的优势。

雍正元年元旦的十一道上谕，我们看到了雍正先事争取地方的策略，为其稳固君权打响了第一砲。此后一连串大量地方事务的工作：摊丁入亩、火耗归公、颁行圣谕广训，乃至豁贱为良、为《钦颁州县事宜》颁发上谕，一方面是刘除允禩集团在地方上盘根错节的力量、继续巩固其君权，一方面则向全国臣民表示，他注重国计民生、他真心关怀爱养百姓，他是一个合格的君王，康熙选择他继位为君是正确的。雍正关怀退休大员、扶持地方宗族的力量，是地方事务的一个环节，也可以放在这样一个脉

\vfill
\begin{BookNotes}
\BookNote{6.}{以上关于康熙晚年诸皇子之间的储位之争，请参见冯尔康老师，《雍正传》(北京：人民出版社，1991年2月第1版第2次印刷），〈第一章储位斗争的胜利者〉，页4-74。}
\BookNote{7.}{参见《大义觉迷录》卷一，叶二十九上，第二道上谕：「…即朕即位以后，而怀藏异志者尚不乏人。说明了这种实情。}
\BookNote{8.}{《大义觉迷录》，卷一，第二道上谕，叶二十六下。}
\end{BookNotes}

% 原书第 340 页（PDF 第 353 页）
\BookPage{340}{353}
\noindent 络里去理解。

因此，雍正正面扶持、利用宗族的力量，不仅止于安靖地方的目的，更有着政治环境的需要。

\BookSubheading{三、乾隆对待宗族的态度：抑制大姓巨室}

\noindent 基本上，如「朱宁三案」，乾隆对于他的父亲雍正让家族取得一部分「处死不法族人」的权力，并没有否决，甚至还继承下来：辛丑（初三日）。谕：刑部参奏，额外主事魏若虚于核议湖北省卢秀等活埋卢在伦身死一案，欲照擅杀例减等问拟，抗不画稿，并自缮奏折说帖，欲行陈奏，实属任性偏执，意图争胜，请交部严加议处一折。朕详加披阅，此案卢秀于缌麻姪孙卢在伦犯窃，设当甫经查获、未交保正之先，因其玷辱祖宗，遽以家法处治致死，犹可援照新例减等。乃卢在伦业经邻人刘惟查获，拏交保正潘见名送官，是此案即应听官办理；而卢秀竟复将卢在伦夺回捆缚，活埋致毙。迨保正查知，欲行禀报，卢秀复阻令勿报寝事。其为畏累逞忿，挟私故杀，自无疑义。该抚照谋故杀本律问拟，情罪原属相当。且为首之卢秀一犯，业已在监病故；所有现拟杖流之卢谷升等，俱系为从之犯，其案情更为无关紧要。今刑部堂官以卢秀因畏卢在伦扳害，起意埋毙，系为自己除害，并非为坫辱祖宗起见，自不得援照擅杀卑幼之例，减等定拟，是以仍照谋故杀本律问拟，所办原属循照向例，并无错误。即核之魏若虚所引四川湖南徐子俸等各案，虽同一尊长杀死卑幼，彼则因玷辱祖宗颜面，忿激致死，故照擅杀律减等；此则因畏其扳害，非激于公忿，故仍照谋故杀问拟。本属两不相悖。乃魏若虚固执己见，任意乖张，始终抗不画稿，且意图见长，并自行缮折，意欲陈奏；其意不过因近来有讦告上司、曾邀迁擢之事，遂欲借此一案，独出己见，希冀一经上闻，或得幸邀识拔，其为进用巧，尤属显然。若部中司员相率效尤，各持异议，堂官何以办事。魏若虚着交部严加议处，以示惩儆。9

\vfill
\begin{BookNotes}
\BookNote{9.}{参见《大清高宗纯皇帝实录》，卷一千四百十二，乾隆五十七年九月辛丑条，关于「刑部参奏额外主事魏若虚于核议湖北省卢秀等活埋卢在伦身死一案」。}
\end{BookNotes}

% 原书第 341 页（PDF 第 354 页）
\BookPage{341}{354}
至于族正的制度，乾隆也延续了下来10。可是乾隆对于宗族的态度有了转变：

\begin{BookQuote}
三十三年，御史张光宪奏请设立大姓族长。谕曰：民间户族繁甚，其中不逞之徒，每因自恃人众，滋生事端；向来聚众械鬬各案，大半起于大姓，乃其明验。惟在地方官实力弹压，有犯必惩，以靖嚣凌之习，政体不过如是。若于各户专立族长名目，无论同宗桀骜子弟未必遽能受其约束；甚者所主非人，必至借端把持，倚强锄弱，重为乡曲之累。正所谓杜弊转以滋弊也。张光宪所请不
\end{BookQuote}

可行。11

大姓族繁势盛，「聚众械斗」的案子，大多数由此而产生，因此乾隆拒绝了御史张光宪奏请于大姓设立族长的建议。而且强调「惟在地方官实力弹压，有犯必惩，以靖嚣凌之习，政体不过如是」，地方事务就应该由地方官来处理；乾隆不愿意宗族、尤其是大姓巨室介入属于地方官权责的事务，其态度是很明显的。

甚至对于族正的功能都有了怀疑，譬如乾隆五十四年福建巡抚徐嗣曾奏请「族正不行阻止举首者，分别治罪。如果教约有方，一岁之中，族内全无命盗械斗等案，给扁奖励；三年无犯，及能将滋事匪徒、查缚送官者，奏给顶戴」，乾隆即给予严厉斥责：

\begin{BookQuote}
但闻闽省民人，俱系聚族而居。该处械斗之风，往往此族与彼族挟嫌争衅，阖姓之人，各行聚众逞凶纠杀，及伤毙人命；即于本族中公议一二人，许以养赡家室，令其顶凶抵命。地方官复因族众难办，又见有凶手承认，遂尔将就完事，屡致酿成大案。今若匪徒犯案，俱责成族中缚送；此等所举族正，大半多系绅衿土豪，未必尽属奉公守法之人。况既族居一处，则滋事者，必有伊子弟亲党在内，族正转得为之包庇，甚而挟嫌妄举，或将衰病者举出充数，滋弊实多。况地方官拘拏人犯，反假手于族正，又给以顶戴，岂不开把持官府之渐。行之日久，将来遇有缉凶拏匪之事，必须向族正索取，竟与世袭土司何异。其法于闽省断非所宜。且此等族正，俱系平民；该省自督抚以下文武员弁，正复不少，遇有民间作奸犯科之事，自应责成该管地方官，认真查办。即或于缉拏
\end{BookQuote}

\vfill
\begin{BookNotes}
\BookNote{10.}{（清）高宗敕撰，《清朝文献通考》，卷十九，〈户口考一·户口丁中赋役〉，「乾隆三十三年」。另外，《大清高宗纯皇帝实录》记载族正的事例甚多，兹不一一列举。}
\BookNote{11.}{（清）高宗敕撰，《清朝文献通考》，卷十九，〈户口考一·户口丁中赋役〉，「乾隆三十三年」。}
\end{BookNotes}

% 原书第 342 页（PDF 第 355 页）
\BookPage{342}{355}
\begin{BookQuote}
凶匪等事，令该处族正举发，亦止可由该处地方官，临时酌办，断无明降谕旨责令专办，并恩给顶戴之理。若概令族正等自行举首，缚送到官，则设地方官何用。是明假以事权，必至倚仗声势，武断乡曲，甚而挟嫌诬首，及顶凶抵命，皆所不免。充其流弊，必至聚众滋事，更复何事不可为。12
\end{BookQuote}

怀疑族正功能的同时，又再度强调地方官员注意该负起的权责。而且乾隆在言论之间，对于聚族而居的大姓巨室，往往「各行聚众逞凶纠杀，及伤毙人命」的情形，很不以为然，也持贬抑的态度；因为此时的大姓巨室，其族长、或族正，已经不再是地方安静的保证，反有可能是「聚族械斗」的领头者、是武断乡曲的绅衿土豪，是地方不能安静的因素。类此情形，《大清高宗纯皇帝实录》所载还有多起事例13，不赘述。

\BookSubheading{四、试析乾隆抑制大姓巨室的原因}

乾隆一如其父、祖，都是乾纲独断、君权独揽的君主；何以对待宗族的态度，如此大异其趣？这要从社会发展的变化说起。

鉴于康熙死前，皇储未定，致使康熙一雍正之间的嬗递并不顺利，雍正继位时的君权不稳固；雍正痛定思痛，制定了储位密建法，使得雍正一乾隆之间的君权顺利转移。因此乾隆继位之时，免除了像雍正当时必须与其他皇亲贵族争权、夺权的情境；可是这并不意味着乾隆的统治若要长治久安，就不需要巩固他的君权。

自顺治元年（1644）满族入关，到康熙二十二年（1683）三藩乱平的四十年间，中原多有战乱，可以不论。自三藩乱平之后，中原内地基本上已无战乱14，直到乾隆中期（乾隆三十年，1765）约共八十二年，以三十年为一代，已可繁衍三代；又值

\vfill
\begin{BookNotes}
\BookNote{12.}{《大清高宗纯皇帝实录》，卷一千三百三十五，乾隆五十四年七月庚戌条。}
\BookNote{13.}{《大清高宗纯皇帝实录》，卷四十九，乾隆二年八月下，「闽浙总督衔专管福建事务郝玉麟议覆」；卷一百三十七，乾隆六年二月下，「广东按察使潘思矩奏」；卷一百九十五，乾隆八年六月下，「闽浙总督那苏图又奏」；卷三百十三，乾隆十三年四月下，「闽浙总督喀尔吉善覆奏」；卷四百八十九，乾隆二十年五月下，「江西巡抚胡宝瑔奏」；卷七百五十九，乾隆三十一年四月下壬戌条，「广东巡抚王检奏」。}
\BookNote{14.}{中原地区无战乱的情形，阎崇年教授在他的书里也有说明。参见阎崇年教授，《康熙大帝》（北京：中华书局，2008年5月1版1刷），页281。}
\end{BookNotes}

% 原书第 343 页（PDF 第 356 页）
\BookPage{343}{356}
\noindent 太平盛世，人口滋生，户口殷繁。因此在乾隆初年以后，因为人口的滋生、积聚，已经出现了不少大姓家族；接近中期以后，强宗豪族更是纷纷涌现。然而强宗豪族的出现，却造成了不少社会问题，武斗如「聚众械斗」，除了上面引述资料外，实录所载，俯拾皆是：

\begin{BookQuote}
闽浙总督衔专管福建事务郝玉麟，议覆署福建陆路提督苏明良奏言，闽省风俗强悍，泉、漳等处尤为好斗生事，请将为首起意、鸣锣聚众之人，从重发遣，
\end{BookQuote}

以靖地方。16

\begin{BookQuote}
闽浙总督那苏图奏报…。又奏，漳浦县吴林二姓，各恃族大丁繁，持械格斗。17壬戌。广东巡抚王检奏，粤民多聚族而居。每族祠置祭田，名为尝租；大户多至数千亩，小户亦有数百亩。租谷按支轮收，除祭祀完粮外，积至盈千累万。赀财丰厚，往往倚强凌弱，恃众暴寡。其势均力敌者，不能取胜，则祠内纠众出斗。议定族中斗伤人，厚给尝租以供药饵；因伤身故者，令木主入祠，给尝租以养妻孥。如伤毙他姓，有肯顶凶认抵者，亦照伤故例。正犯漏网奸徒愈无
\end{BookQuote}

顾忌。18

至于文斗则有「聚讼」，例如：

\begin{BookQuote}
广东按察使潘思矩奏，粤民多聚族而居，各建宗祠。置尝租，岁入，实费于祭祀及给族人等用，甚鲜。余以生息，月积岁累。偶与外姓睚眦小忿，通族扛帮争讼，一切费用，取给尝租。19
\end{BookQuote}

此类情形则以乾隆二十九年江西巡抚辅德的奏报最具代表性：

\begin{BookQuote}
奏为查禁祠宇流弊，涤讼源，以维风教事。窃照江西民情健讼，有司勤惰不齐，州县自理词讼及上司批查案件，多不遵照例限审结，且有判断失平，不能
\end{BookQuote}

\vfill
\begin{BookNotes}
\BookNote{15.}{有关乾隆时期人口滋长的问题，当时已有人论述，其中最具代表性的是洪亮吉的说法。洪亮吉有关「人口论」的论述，请参考：吴德玲着，《洪亮吉《意言》研究》（台北，花木兰文化出版社，2006年9月；收入潘美月、杜洁祥主编，《古典文献研究辑刊》三编，第二八册），尤其第六章，页87-119。}
\BookNote{16.}{《大清高宗纯皇帝实录》卷四十九，乾隆二年八月下，「闽浙总督衔专管福建事务郝玉麟议覆」。}
\BookNote{17.}{《大清高宗纯皇帝实录》卷一百九十五，乾隆八年六月下，「闽浙总督那苏图又奏」。}
\BookNote{18.}{《大清高宗纯皇帝实录》卷七百五十九，乾隆三十一年四月下，「广东巡抚王检奏」。}
\BookNote{19.}{《大清高宗纯皇帝实录》卷一百三十七，乾隆六年二月下，「广东按察使潘思矩奏」。}
\end{BookNotes}

% 原书第 344 页（PDF 第 357 页）
\BookPage{344}{357}
折服其心，未免益长刁风，而滋拖累。臣到任以来，逐一清查…惟查各属讼案繁多之故，缘江西民人有合族建祠之习，本籍城乡暨其郡郭并省会地方，但系同府同省之同姓，即纠歛银钱，修建祠堂，率皆栋宇辉煌，规模宏敞，其用余银两置产收租，因而不肖之徒，从中觊觎，每以风影之事，妄启讼端，藉称合族公事，开销祠费。县讼不胜，即赴府翻；府审批诘，又赴省控。何处控诉，即住何处祠堂，即用何处祠费；用竣复按户派出私财，任意侵用。是祠堂有费，实为健讼之资；同姓立祠，竟为聚讼之地，欲弭讼端，不可不清其源而塞其流也。臣查民间祠堂，如系建于本乡，时祭飨而联络族谊，设公费以教养子弟，乃系敦尚古道，实为美俗可风。若远建于府省地方，祭飨无闻，族谊不浃，其屋宇则傍宿健讼之徒，其公财则积为逞讼之费，颓风败习，莫此为甚。…原其创建之初，不过一二好事之徒，借端建议，希图经手侵渔，访其同府同省同姓，或联络于生童应考之时，或偏走之农民收割之后，百计劝捐，多方耸动，愚民溺于习俗，乐为输助，故其费日集而多，其风日踵而盛。初成广厦，置之空闲，歇讼聚赌，窝匪藏奸，不可究诘。近于省会祠中，复经拏获私铸案犯。…况滥觞至为聚讼之薮，甚而窝赌窝匪，无所不有，尤难听其流弊无所底止。…倘准留之房，仍有讼棍盘踞，及窝赌窝匪情事，除严拏本犯治罪外，即将其屋宇入官，或作堆铺，或给未建衙署之员弁居住。…嗣后永远不许添建府省祠堂公宇，其有实系敦本支而睦宗族者，祇许于本乡本村以时飨祀，庶几礼教可明，讼源可涤，而民生日厚矣。20

但是无论是武斗，或是文斗，同样都会造成社会治安的问题。这种情形，任何一个有为的君主都是不能容忍的，何况是乾纲独断、君权独揽的乾隆皇帝。乾隆为了民生的安定、为了巩固他的统治与君权，因此而存有抑制大姓巨室的态度是可以想见的。

\BookSubheading{结语}

\noindent 从以上的叙述可知，雍正与乾隆父子二人，在对待宗族的态度上，有相同处，也有相异之处。

\vfill
\begin{BookNotes}
\BookNote{20.}{辅德，〈奏为查禁祠宇流弊，涤讼源，以维风教事〉，国立故宫博物院藏，文献编号403017412。}
\end{BookNotes}

% 原书第 345 页（PDF 第 358 页）
\BookPage{345}{358}
相同之处是，因为宗族向来都是地方上一股重要的力量，为了安靖地方，也为了稳固君权与长治久安的需要，雍正乾隆父子二人都必须利用宗族。因此雍正不只设置了族正，以协助地方保甲力所不及之处，甚至还给予宗族处死不法族人的权力。而乾隆在这方面，基本上并没有推翻雍正的做法，一方面仍然延续了族正的制度，另一方面也没有完全否决雍正时让家族取得一部分「处死不法族人」的权力；不过，必须注意的是，乾隆的态度与作法有了改变。

不同之处是，雍正继位之时，政敌环伺，君权是不稳固的；为了稳固君权，他从地方着手，争取地方官与全国百姓的支持，因此他做了大量的地方事务；宗族既然是地方上的重要组成，而雍正初年似乎尚未显现其力量，于是雍正正面扶持宗族，给予宗族若干行政上的权力。显然雍正的作为有着政治环境的需要。乾隆继位之后，不像其父面临君权的威胁，但社会情境却有了巨大的变化，那就是人口的滋长，导致大姓巨室的出现；原来期待利用宗族成为一股安靖社会的力量，大姓巨室的出现，不断发生宗族之间的械斗，反成为社会不能安定的力量；因此乾隆继承、延续雍正政策的同时，他并没有僵固不化，一方面利用宗族，另一方面则反过来抑制大姓巨室，以免侵蚀了他的君权与统治。在此我们看到，乾隆灵活的统治手腕，适应了社会新形势的变化

总之，雍正乾隆父子在巩固君权与统治的同一理念下，他们对待宗族的态度与作法，虽然因为对应社会情境的变化而有所差异，他们的目的一地方的安静，君权的稳固，统治的顺畅一是相同的。这种差异所获致的结果，反倒是殊途同归的。


% 原书第 346 页（PDF 第 359 页）
\BookPage{346}{359}
